\chapter{Systematic Literature Review}
\label{app:systematic-literature-review}
\automark[section]{chapter}
\markboth{Systematic Literature Review}{Systematic Literature Review}

\section{Purpose and Review Question}

The review examines how operational telemetry supports LLM-based fault
diagnosis, with particular attention to the evidence used to evaluate
correctness and reliability. Its purpose is to derive criteria for choosing
and evaluating an agent approach for the Norlys PoC. It does not assume
that the highest published score identifies the best transferable system.
The review question is:

\begin{quote}
How do LLM-based agents use operational telemetry to diagnose faults in
cloud and microservice systems, and what evidence supports the correctness
and reliability of their diagnoses?
\end{quote}

The synthesis in Chapter~\ref{cha:state-of-the-art} connects the findings
to design and evaluation choices. This appendix documents the method,
selection outcome, appraisal and retained reports. The assessment workbook
contains the record-level decisions and detailed extraction.

The review question concerns diagnosis. It provides the systematic
evidence for diagnostic evaluation (RQ3) and part of the evidence
requirements for telemetry correlation (RQ1). Security telemetry (RQ2) and
remediation (RQ4) were not reviewed systematically. The search captured
some reports on these topics, but remediation-only and attack-investigation
reports were excluded under E3 and E2. Chapter~\ref{cha:state-of-the-art}
discusses them as context sources.

Context sources are cited for background and contrast, not counted as
review evidence. Six of them were returned by the searches and excluded
at title/abstract screening under E3 because they do not use an LLM for
diagnosis. These are Sage (SLR0679), Eadro (SLR0015), CIRCA (SLR0609),
TraceWeaver (SLR0681), SIERRA (SLR0206) and the causal-inference evaluation
by Pham et~al.\ (SLR0146). Holmes and the dataset study by Engelen
et~al.\ were not returned by the searches. Holmes investigates DDoS attacks
and would fall under E2. PACE-LM, which Chapter~\ref{cha:state-of-the-art}
cites for confidence calibration, is not a context source. It is the
assessed version of the retained report LM-PACE (SLR0613).

\section{Sources and Search Strategy}

The main searches used IEEE Xplore and the ACM Digital Library Full-Text
Collection on 16 October 2026. Google Scholar supplemented these exports
on 22 October 2026 with three focused queries. Two previously identified
benchmark reports, OpenRCA and ITBench, were added separately as citation
additions. This was not an exhaustive citation-network search.

The publication-year range was 2020--2026, with English-language reports
eligible for screening. IEEE used conference, journal and early-access
content types. The ACM export used the e-publication date range
1 January 2020--31 December 2026. Coverage ends at the respective search
dates. Additional ACM interface filters were not recorded and are not
assumed.

The database queries deliberately covered operational diagnosis more
broadly than LLM agents. The narrower LLM-focused question was applied
during selection. The LLM restriction was a screening criterion, so a
database result did not need to mention an LLM in its abstract to be
captured. This made the database exports sensitive but not specific. Most
of the 793 database candidates describe statistical, causal or
deep-learning diagnosis without an LLM, and most title/abstract exclusions
are E3 for that reason. Google Scholar supplemented these records with
more focused LLM and agent queries and with publishers beyond the two main
collections. It was not treated as a complete or stable enumeration of the
field.

\subsection{Database Expressions}

The following is the recorded ACM expression, applied to the abstract
field. Wildcards and alternatives are preserved from the search record.

\begin{quote}\small\raggedright
\texttt{(cloud OR microservice* OR "distributed systems" OR "online service" OR SIEM)}\\[3pt]
\texttt{AND ("root cause" OR "fault localization" OR "fault localisation" OR troubleshooting OR "incident diagnosis" OR "distributed tracing" OR remediation OR "self-healing" OR ((security OR attack) AND investigation))}\\[3pt]
\texttt{AND (log OR logs OR trace* OR tracing OR metric* OR telemetry OR "security events" OR "large language model" OR "large language models" OR LLM OR causal)}
\end{quote}

IEEE used the same three concept blocks. Each term was explicitly
qualified with the abstract field. Its recorded localisation term was
\texttt{"fault locali?ation"}, and its language-model term was
\texttt{"large language model*"}. The exact executable expression, rather
than a retrospectively improved query, is preserved in the workbook's
Searches sheet. A copy is held in \texttt{Review/searches.csv}. The two
expressions differ in one respect. ACM quotes \texttt{"self-healing"} as a
phrase, while IEEE leaves the hyphenated term unquoted. Its effect on the
IEEE result set was not tested. The difference is recorded rather than
corrected because the exports were made with these expressions.

\subsection{Google Scholar Supplement}

The three queries were:
\begin{enumerate}
\item \texttt{"root cause analysis" "agent" "cloud" "language"}
\item \texttt{"microservices" "diagnosis" "large language model"}
\item \texttt{"benchmark" "cloud" "agents" "diagnosis"}
\end{enumerate}

For each query, the first 50 relevance-ranked results were captured using
the 2020--2026 range, with patents disabled and citation records enabled.
This produced 150 occurrences. The limit applies to captured records,
not to the reported size of Scholar's result set. Relevance ordering and
the bounded capture limit restrict recall and reproducibility.

\section{Eligibility and Selection}

A report was eligible when it presented original research on diagnostic
use of operational telemetry in cloud services, microservices or supporting
infrastructure, with an LLM-based diagnostic contribution. Diagnostic
benchmarks were also eligible when they provided usable tasks, reference
causes and scoring procedures for evaluating such agents. Such a benchmark
qualified even when the report itself evaluated no LLM. Diagnosis in
networks, databases and GPU clusters was treated as supporting
infrastructure when the diagnosed system carried or hosted cloud services,
such as data-centre, cloud and mobile-core networks. Optical
transport-layer control was treated as outside this scope. A relevant study
could use logs, metrics, traces, incident records or combinations of these.
It did not have to use the PoC's particular vendors or agent organisation.

Research preprints were eligible, with the retrieved version recorded.
They were not treated as equivalent to independently peer-reviewed
publications. Earlier and later reports of the same work were linked when
their contributions overlapped. Overlapping evaluations were not counted
as independent confirmations. Following the Cochrane distinction between studies and
reports, reports were linked as one study only when they shared an
experiment, a dataset or results, and not on author overlap alone
\parencite{li2024cochranech5}. VeriTrail-RCA (SLR0057) and VeriOps-Agent
(SLR0266) share four authors and appeared in the same year. Their
evaluation sets, primary outcomes and results differ, and neither cites
the other. The first evaluates root cause localisation on public benchmark
subsets, while the second evaluates remediation gating on Kubernetes fault
scenarios. They are therefore retained as separate studies, but as work
from one group they are not treated as independent confirmations of
evidence verification.

\begin{table}[htbp]
\centering\small
\caption{Exclusion criteria applied during selection.}
\label{tab:slr:exclusions}
\begin{tabularx}{\textwidth}{@{}lX@{}}
\toprule
Code & Reason \\
\midrule
E1 & Ineligible publication type: review, dissertation, proceedings volume or proposal without completed original results. \\
E2 & Outside cloud/service operations and supporting infrastructure. \\
E3 & No LLM-based diagnostic contribution and no usable diagnostic benchmark. Detection alone, code repair and remediation-only evaluation are outside scope. \\
E4 & Insufficient completed method or assessable evaluation after examining the full text. \\
E5 & Retracted publication or retraction notice. \\
\bottomrule
\end{tabularx}
\end{table}

E1 takes precedence for reports that only propose future work. E4 applies
when a report claims a completed method or experiment but does not describe
it sufficiently for appraisal. A position-paper label alone does not
establish ineligibility.

A weak result or an unfavourable comparison was not an exclusion reason.
E4 distinguishes the absence of an assessable completed evaluation from
weak evidence in an otherwise identifiable evaluation. For example,
SuperDRI describes an implemented workflow, held-out incident comparison
and query-execution counts. It is retained for its method and process
measurements, while its lack of diagnostic correctness evidence limits
what can be concluded from it. GS2-14 is excluded under E4 because its
reported accuracy lacks a fault-case
count, reference-label procedure and sufficiently described baselines.
That borderline assessment remains identifiable in the workbook.

\subsection{Duplicate Removal and Screening}

The database exports contained 819 occurrences: 550 IEEE and 269 ACM.
Duplicate matching used DOI equality, exact normalised titles and a checked
title variant. The recorded 26 duplicate occurrences comprised 19 DOI
matches, six exact-title matches and one title variant, leaving 793
candidates. Scholar results were checked against these records and earlier
Scholar captures, removing a further 46 repeated occurrences. Titles and
metadata were used to distinguish repeats from different reports of the
same study. The latter were linked as companions where appropriate.

Title and abstract screening assessed relevance against the question and
criteria. Records retained at this stage proceeded to full-text retrieval.
Retention did not itself constitute inclusion. The 702 title/abstract
exclusions comprise 47 E1, 22 E2, 631 E3 and two E5 decisions. Most E3
decisions concern database records on non-LLM diagnosis, and these share a
standard reason for the same criterion. Full-text assessments examined the
method, diagnostic task, evaluation and relevance before assigning
inclusion, exclusion or companion status. Each full-text decision has a
record-specific reason. Missing full text is a separate retrieval outcome,
not evidence that a paper fails the relevance criteria.

I checked all 113 full-text decisions and 176 title/abstract decisions,
including the borderline cases, and retained the original decisions.
My decisions and reasons are recorded in the assessment workbook.
I used GitHub Copilot to suggest title/abstract decisions, locate open
full texts, prepare extraction fields and check reported numbers against
the source passages.

\subsection{Method and Amendments}

The review separates search, selection, extraction and appraisal, following
the methodological structure described by Kitchenham and Charters
\parencite{kitchenham2007guidelines}. Search reporting records expressions,
sources, dates and limits in line with PRISMA-S
\parencite{rethlefsen2021prismas}.

The review had no registered protocol and no screening pilot. Three
deviations from a prospective protocol are recorded. First, the
LLM-focused question was agreed after the broad database searches and
applied during screening, so the database expressions do not require an
LLM term. Second, full-text retrieval was repeated. The first pass assessed
87 of the 197 reports sought. Renewed passes on 29 and 30 October 2026
searched publisher open-access versions, author manuscripts, institutional
repositories and preprint servers for the remainder. The first renewed pass
added 13 assessments, with four inclusions and nine exclusions. The second
added 13 more, with eight inclusions and five exclusions. Third, exclusion
codes were harmonised. Proposal-only reports were recoded to E1, and one
title/abstract E4 decision (SLR0570) was recoded to E3 because E4 applies
only after full-text examination. The search
expressions and candidate set were unchanged throughout, and earlier
decisions and retrieval attempts remain in the review records.

\subsection{Selection Outcome}

\begin{table}[htbp]
\centering\small
\caption{Reconciled selection counts. Occurrences, reports and independent
findings are not interchangeable units.}
\label{tab:slr:flow}
\begin{tabularx}{\textwidth}{@{}Xr@{}}
\toprule
Stage & Count \\
\midrule
IEEE and ACM occurrences & 819 \\
Database duplicate occurrences removed & 26 \\
Unique database candidates & 793 \\
Scholar occurrences captured & 150 \\
Scholar duplicates removed & 46 \\
New Scholar candidates & 104 \\
Separately recorded citation additions & 2 \\
Candidates after deduplication and additions & 899 \\
Title/abstract exclusions (667 database, 35 Scholar) & 702 \\
Reports sought for full-text assessment & 197 \\
Reports not retrieved & 84 \\
Reports assessed in full text & 113 \\
Full-text exclusions & 34 \\
Companion reports linked to retained work & 3 \\
Reports retained for synthesis & 76 \\
\bottomrule
\end{tabularx}
\end{table}

The full-text exclusions comprise four E1, three E2, twenty-two E3 and five
E4 decisions. The three companions link the earlier AIOpsLab report
(SLR0553) to GS3-04, the earlier recursion-of-thought report (GS1-35) to
SLR0071, and the earlier GALA report (GS1-40) to SLR0600. They are recorded
without adding independent evidence to the same study families.

Of 197 reports sought, 84 were not retrieved. Retrieval used open routes:
publisher open-access versions, author manuscripts, institutional
repositories and preprint servers, searched by DOI and exact title.
Subscription access through the university library was not used. For 55
of the 84 reports, the only version located was subscription-only. Fifty
of these are in IEEE Xplore, three at Springer and two in the ACM Digital
Library. For 19 reports an open copy was located, but the host's automated
bot check prevented download during the retrieval passes. One report was
found only as an author upload to ResearchGate, and nine had no full-text
link and no matching version could be located. The total comprises 62
database records and 22 Scholar records. Following PRISMA 2020, these
reports are recorded as not retrieved rather than excluded
\parencite{page2021prisma}. The record of each attempt is kept in
\texttt{Review/retrieval-log.csv}.
Retrieval was closed after the second renewed pass. The unretrieved reports
are treated as a declared limitation (Section~\ref{sec:slr:limitations}),
and the sensitivity check below estimates their possible effect.

A sensitivity check estimated how the unretrieved reports could affect
the findings. Each was rated from its abstract, or its title where no
abstract was available, against the inclusion criteria. Thirty-six appear
likely to be eligible, 31 are unclear and 17 appear likely to be excluded.
The likely eligible and unclear reports resemble the retained set, with
fewer benchmarks. They mostly describe fixed LLM workflows, multi-agent
systems and graph-guided investigation, and 11 claim some form of evidence
verification. Abstracts
rarely state how the evaluation was controlled, so the check cannot show
whether these reports would change the appraisal counts. The missing
reports are therefore a coverage limitation for design approaches,
particularly evidence verification, rather than evidence of a different
evaluation practice.

The supplementary sources added substantially to the review. Twenty-eight
of the 76 retained reports (37\,\%) were found only through the Scholar
supplement or the citation additions. Sixteen of these are preprints. The
remainder appeared at machine-learning venues such as AAAI, EMNLP, ICLR,
MLSys and NeurIPS, or with publishers outside the two collections, such as
Springer and MDPI. They include 12 of the 18 benchmark and framework
reports. IEEE Xplore and the ACM Digital Library alone would have missed
most of the benchmark evidence.

\section{Extraction and Quality Appraisal}

For retained reports, extraction records the diagnostic method, system
context, telemetry, data and reference causes, evaluation procedure,
reported findings, limitations and implications for the PoC. Source URLs,
retrieved versions and passage locations support checking the extraction.
The workbook's detailed fields additionally distinguish agent roles,
models, tools, baselines, metric definitions, efficiency and reliability.

Appraisal examines the following dimensions rather than assigning a total
quality score:
\begin{enumerate}
\item \textbf{Task and ground truth:} what constitutes a correct diagnosis,
how faults and labels are established, and which cases are omitted.
\item \textbf{Comparison validity:} shared evidence, models, training access,
tool interfaces and inference budgets, including whether baselines were actually rerun.
\item \textbf{Outcome validity:} separation of detection, localisation,
causal explanation, tool success, task completion and recovery.
\item \textbf{Reliability:} repeated runs, healthy controls, missing or
misleading evidence, uncertainty, timeouts and incomplete answers.
\item \textbf{Transfer and reproducibility:} system realism, proprietary
inputs, released artifacts, model versions and reported operating cost.
\end{enumerate}

The synthesis groups studies by these evaluation and design issues. It
uses quantitative results only with their metric and experimental context.
There is no pooled accuracy or league table across different datasets.
For example, a successful KQL call, a penalised root-cause score and a
joint cause-and-remedy success rate cannot be treated as the same outcome.
The number of retained reports is an outcome of selection, not a target
quota.

\subsection{Appraisal Results}
\label{sec:slr:appraisal-results}

Each dimension was rated strong, partial or weak for every retained report.
A strong rating means the report states the information and the design
supports the claim it makes. A partial rating means the information is
present but incomplete or confounded. Table~\ref{tab:slr:appraisal}
summarises the ratings for the 67 reports appraised against their full
text. Nine further reports (SLR0009, SLR0048, SLR0057, SLR0167, SLR0183,
SLR0256, SLR0266, SLR0328 and SLR0563) were appraised from their
extraction records. Their eligibility and extraction were based on the
full text, but the author-uploaded copies could not be downloaded again
when the appraisal was made. Fields that the extraction records did not
cover were left unrated, and these nine reports are not counted in
Table~\ref{tab:slr:appraisal}.

\begin{table}[htbp]
\centering\small
\caption{Appraisal ratings for the 67 retained reports appraised against
their full text.}
\label{tab:slr:appraisal}
\begin{tabularx}{\textwidth}{@{}Xrrr@{}}
\toprule
Dimension & Strong & Partial & Weak \\
\midrule
Task and ground truth & 23 & 38 & 6 \\
Comparison validity & 6 & 46 & 15 \\
Outcome validity & 18 & 41 & 8 \\
Reliability & 4 & 29 & 34 \\
Transfer and reproducibility & 9 & 27 & 31 \\
\bottomrule
\end{tabularx}
\end{table}

Reliability is the weakest dimension. Of the 67 reports, 20 repeat runs
and 17 report the variability between them. Ten include healthy or
no-fault cases, and seven let the system report that the evidence is
insufficient. Eighteen test missing, noisy or misleading evidence.
Fourteen check whether a diagnosis contains claims that the evidence does
not support.

Comparison validity is mostly partial. Seven reports match the inference
budget across the compared methods. Twenty-five report cost or steps
without matching them, and 29 report neither. Thirty-five compare only
against other LLM methods, and seven have no baseline. In 27 reports the
proposed method uses training or historical data that its baselines do not,
so a difference in score often reflects more than the design being tested.

Transfer is limited by data access. Eighteen reports evaluate only
proprietary incidents and thirteen combine proprietary and public data.
Twenty-three release code or data. Sixteen report time, tokens and cost
in full, and 34 report some of these.

Outcome measures vary. Twenty-four reports score top-$k$ localisation,
fourteen require an exact or joint cause, and twelve report classification
accuracy or F1 over fault types. Six use an LLM judge, five use textual
similarity to a reference and four report task success. Six of the 67 do
not score diagnostic correctness against a ground truth at all. These
counts support the evaluation criteria in
Section~\ref{sec:state-of-the-art:evaluation}.

\section{Limitations}
\label{sec:slr:limitations}

Coverage is limited by the selected collections, search expressions,
publication window and bounded Scholar supplement. The 84 unretrieved
reports are a substantial limitation. The assessed set cannot show that no
other relevant approach exists, and it cannot identify a best-performing
system. Retrieval through open routes also favours preprints and
open-access publishers. Fifty of the unretrieved reports are
subscription-only IEEE papers, so IEEE conference and journal work is
under-represented among the retained reports.

Nine retained reports were appraised from their extraction records rather
than from the full text, so their appraisal is less complete.

There was no second independent reviewer, which leaves a risk of selection
and interpretation errors. Preprints may change after the retrieved version.
Several studies use proprietary incidents, small fault sets or
model-dependent judges. These limitations affect how confidently the
findings can be applied to the PoC. The review provides evaluation
hypotheses and design criteria whose applicability must be tested in that
environment.

\section{Retained Reports and Appraisal Summary}

Table~\ref{tab:slr:catalogue} identifies all 76 retained reports using the
stable screening IDs. Titles link to the assessed source version. The
brief appraisal highlights a limitation rather than reproducing the full
extraction. Detailed methods, findings and passage references remain in
\texttt{SLR-screening-completed.xlsx}, especially the Studies, Appraisal and
Key results sheets.

\begingroup
\small
\setlength{\tabcolsep}{3pt}
\renewcommand{\arraystretch}{1.12}
\begin{longtable}{@{}p{0.14\textwidth}p{0.43\textwidth}p{0.39\textwidth}@{}}
\caption{Retained reports, approach and principal appraisal concern.}
\label{tab:slr:catalogue}\\
\toprule ID & Report and approach & Appraisal concern \\
\midrule\endfirsthead
\multicolumn{3}{l}{\textit{Table \thetable\ continued}}\\
\toprule ID & Report and approach & Appraisal concern \\
\midrule\endhead
\midrule\multicolumn{3}{r}{\textit{Continued on next page}}\\\endfoot
\bottomrule\endlastfoot
SLR0057 & \href{https://www.researchgate.net/publication/406279260_VeriTrail-RCA_Evidence-Verified_Root_Cause_Localization_for_Microservice_Backends}{VeriTrail-RCA: Evidence-Verified Root Cause Localization for Microservice Backends}\par\emph{Evidence-verifying diagnostic agent} & Only one backbone. Subset selection, inconsistent descriptions of some aggregate/process denominators and unreleased sampling code reduce confidence. \\
SLR0059 & \href{https://arxiv.org/pdf/2605.15611}{TopoEvo: A Topology-Aware Self-Evolving Multi-Agent Framework for Root Cause Analysis in Microservices}\par\emph{Topology-guided multi-agent diagnosis} & Incident counts, split independence and repeated-run uncertainty are not sufficiently specified in the evaluation. Topology faithfulness partly uses the method’s own salient edges. \\
SLR0071 & \href{https://xplorestaging.ieee.org/ielx8/8858/11676280/11575089.pdf?arnumber=11575089}{Toward In-Depth Root Cause Localization for Microservices With Multi-Agent Recursion-of-Thought}\par\emph{Recursive multi-agent localisation} & The improved result is not universal across datasets. Aggregate localisation does not validate every explanation. \\
SLR0082 & \href{https://arxiv.org/pdf/2603.21145}{NeSy-Edge: Neuro-Symbolic Trustworthy Self-Healing in the Computing Continuum}\par\emph{Neuro-symbolic agentic diagnosis} & Log-centric constructed evaluation. Independence and verification of derived root-cause annotations are not fully established. \\
SLR0096 & \href{https://arxiv.org/pdf/2509.26463}{ErrorPrism: Reconstructing Error Propagation Paths in Cloud Service Systems}\par\emph{Code-assisted diagnostic agent} & The static-analysis baseline is scored by candidate-function precision, not the same exact-path metric. Private Go systems and selected templates limit transfer. \\
SLR0102 & \href{https://arxiv.org/pdf/2605.15581}{STAR: A Stage-Attributed Triage and Repair Framework for RCA Agents in Microservices}\par\emph{Agent trajectory correction} & Reimplemented workflows, unclear incident denominators and judge-dependent stage labels. Extra inference budget is part of the improvement. \\
SLR0119 & \href{https://arxiv.org/pdf/2602.05292}{ORACL: Optimized Reasoning for Autoscaling via Chain of Thought With LLMs for Microservices}\par\emph{Trained diagnostic and control workflow} & RCA sample denominator, class aggregation and split details are not sufficiently clear in the inspected evaluation. No RCA confidence interval. \\
SLR0131 & \href{https://arxiv.org/pdf/2509.25874}{LogPilot: Intent-aware and Scalable Alert Diagnosis for Large-scale Online Service Systems}\par\emph{Alert-guided log diagnosis} & Private selected alerts and subjective usefulness. No conventional baseline is tested. \\
SLR0171 & \href{https://arxiv.org/pdf/2506.10043}{TrioXpert: An Automated Incident Management Framework for Microservice System}\par\emph{Modality-specialised multi-agent diagnosis} & Incident timestamps are supplied for classification/localisation. Private-domain transfer, unclear split leakage and limited uncertainty reporting. \\
SLR0208 & \href{https://arxiv.org/pdf/2507.03224}{RCA Copilot: Transforming Network Data into Actionable Insights via Large Language Models}\par\emph{Retrieval-assisted network diagnosis} & Very small dataset, limited held-out topology evidence and no exact component/cause correctness metric. Semantic similarity is not diagnostic accuracy. \\
SLR0245 & \href{https://xplorestaging.ieee.org/ielx8/6287639/11323511/11598804.pdf?arnumber=11598804}{CausalOpsBench: A Reproducible Interventional-Replay Benchmark for Topology-Aware Operational Incident Response}\par\emph{Interventional diagnostic benchmark} & Mostly synthetic cross-domain scaffold with compact candidate labels. External and native composites cannot be compared directly. \\
SLR0256 & \href{https://www.preprints.org/frontend/manuscript/41c58f87035172007cb15106de9bf93a/download_pub}{Hierarchical Expert Multi-Agent Framework for Causal Root Cause Localization in Cloud-Native Microservices}\par\emph{Hierarchical expert diagnosis} & Preprint with sparse reproducibility, ground-truth and split details. Narrative names a best baseline absent from its comparison table. \\
SLR0303 & \href{https://arxiv.org/pdf/2503.23051}{COCA: Generative Root Cause Analysis for Distributed Systems with Code Knowledge}\par\emph{Code-assisted generative diagnosis} & Public resolved issues risk pretraining contamination. Ground-truth creators also rate usefulness. \\
SLR0304 & \href{https://doi.org/10.1109/tsc.2025.3629066}{TAMO:Fine-Grained Root Cause Analysis via Tool-Assisted LLM Agent With Multi-Modality Observation Data in Cloud-Native Systems}\par\emph{Hybrid ML tools with LLM interpretation} & Most quantitative evidence evaluates ML tools, not final LLM correctness. Reported subsecond inference is not demonstrated as complete LLM investigation time. \\
SLR0359 & \href{https://arxiv.org/pdf/2301.03797}{Recommending Root-Cause and Mitigation Steps for Cloud Incidents Using Large Language Models}\par\emph{Fine-tuned incident recommendation} & Private historical descriptions, not autonomous live investigation. Humans select the best of several LLM candidates while encoder baselines supply one, limiting comparison. \\
SLR0571 & \href{https://arxiv.org/pdf/2603.02032}{MetaRCA: A Generalizable Root Cause Analysis Framework for Cloud-Native Systems Powered by Meta Causal Knowledge}\par\emph{Offline LLM knowledge with statistical diagnosis} & Offline knowledge construction is excluded from online latency. Eighteen cases with missing trace relationships are excluded. \\
SLR0572 & \href{https://arxiv.org/pdf/2601.02736}{Hypothesize-Then-Verify: Speculative Root Cause Analysis for Microservices with Pathwise Parallelism}\par\emph{Parallel hypothesis verification} & Different backbones and training access confound the architecture comparison. Brief preliminary study, with limited sample/split and uncertainty detail. \\
SLR0573 & \href{https://arxiv.org/pdf/2401.13810}{Automated Root Causing of Cloud Incidents using In-Context Learning with GPT-4}\par\emph{Retrieval-based incident recommendation} & Private historical descriptions and limited human sample. Model size and context configuration change with the compared approach. \\
SLR0576 & \href{https://arxiv.org/pdf/2601.22208}{Stalled, Biased, and Confused: Uncovering Reasoning Failures in LLMs for Cloud-Based Root Cause Analysis}\par\emph{Controlled agent evaluation} & Graph-mediated observations and selected models limit transfer. Valid graph paths are not proof of causality. \\
SLR0577 & \href{https://nkcs.iops.ai/wp-content/uploads/2025/03/ChangeLLM.pdf}{A Multimodal Intelligent Change Assessment Framework for Microservice Systems Based on Large Language Models}\par\emph{LLM-assisted change diagnosis} & Selected private cases and supplied fault times. Main table reports D1 Avg@5 0.975 with Top1 0.775, inconsistent with the stated mean-of-top-k formula if all accuracies are bounded by one. \\
SLR0578 & \href{https://saurabhjha.one/pubs/ASE2022/paper.pdf}{A fault injection platform for learning AIOps models}\par\emph{Diagnostic evaluation infrastructure} & Retained as evaluation infrastructure with injected-fault ground truth. No LLM is evaluated and no task set is released, so this is the weakest fit to the benchmark criterion. Two benchmark applications and private infrastructure. \\
SLR0579 & \href{https://arxiv.org/pdf/2603.02057}{Beyond Microservices: Testing Web-Scale RCA Methods on GPU-Driven LLM Workloads}\par\emph{Diagnostic transfer benchmark} & Small topology, single faults and oracle timestamps. An LLM is the monitored workload, not the diagnostician. \\
SLR0586 & \href{https://alessandrocornacchia.github.io/assets/pdf/apsys25-llm-failure-analysis-cornacchia.pdf}{Between Promise and Pain: The Reality of Automating Failure Analysis in Microservices with LLMs}\par\emph{Tool and agent comparison} & Very small diagnostic problem set and different reasoning models. Code-QA results are not incident diagnosis scores. \\
SLR0593 & \href{https://arxiv.org/pdf/2305.15778}{Automatic Root Cause Analysis via Large Language Models for Cloud Incidents}\par\emph{Retrieval-assisted diagnostic workflow} & One private service and pre-existing handler expertise. Temporal independence of the split is not established. \\
SLR0595 & \href{https://arxiv.org/pdf/2607.01788}{KRCA: An Efficient Root Cause Analysis System in Hyper-scale Microservice Systems via Agentic AI}\par\emph{Graph and memory guided agent diagnosis} & Private hyper-scale infrastructure and historical memory. Mixed-model main/subagents and absence of conventional baselines prevent isolating pure architecture superiority. \\
SLR0600 & \href{https://arxiv.org/pdf/2608.08968}{GALA: Graph-Augmented LLM Agents for Root Cause Analysis and Incident Response in Microservices}\par\emph{Graph-guided multi-agent workflow} & Pod match does not require correct fault type. Report rubric partly rewards remediation absent from some baselines. \\
SLR0602 & \href{https://arxiv.org/pdf/2504.11505}{eARCO: Efficient Automated Root Cause Analysis with Prompt Optimization}\par\emph{Prompt-optimised incident diagnosis} & Private data, summarised references and model-judge dependence. Reported gains are ordinal rating improvements, not exact diagnosis accuracy. \\
SLR0603 & \href{https://grace-liu.github.io/static/papers/25-SIGCOMM-bian.pdf}{Towards LLM-Based Failure Localization in Production-Scale Networks}\par\emph{Production network localisation} & Monitoring first restricts candidates to six devices with reported >98\% coverage. Private network and device localisation are narrower than open-ended root-cause explanation. \\
SLR0612 & \href{https://arxiv.org/pdf/2502.08224}{Flow-of-Action: SOP Enhanced LLM-Based Multi-Agent System for Root Cause Analysis}\par\emph{Runbook-guided multi-agent diagnosis} & Coverage scores penalise false causes only by 0.1 and are not exact accuracy. Baselines differ in accessible telemetry. \\
SLR0617 & \href{https://arxiv.org/pdf/2404.03662}{X-Lifecycle Learning for Cloud Incident Management using LLMs}\par\emph{Dependency-context incident diagnosis} & Private selected service and historical records. Fine-tuned GPT-3 versus GPT-4 comparisons confound training method and backbone. \\
SLR0618 & \href{https://arxiv.org/pdf/2601.02732}{Agentic Memory Enhanced Recursive Reasoning for Root Cause Localization in Microservices}\par\emph{Tool-using agent with memory} & Incident counts and repeat uncertainty are insufficiently clear. Per-alert versus batch timing differs across baselines. \\
SLR0619 & \href{https://arxiv.org/pdf/2310.16340}{RCAgent: Cloud Root Cause Analysis by Autonomous Agents with Tool-Augmented Large Language Models}\par\emph{Tool-using diagnostic agent} & Private single-platform data and selection through an existing diagnostic tool. LLM grading and human helpfulness do not establish exact-cause accuracy. \\
SLR0623 & \href{https://nkcs.iops.ai/wp-content/uploads/2026/04/Yujia__Aloha_to_FSE_26.pdf}{Aloha: Localizing Batch Failures in Large-scale Cloud Systems via Contrast Analysis and Human-in-the-Loop Agent}\par\emph{Agent-assisted contrast analysis} & Private cases, limited label and repetition detail. Human corrections contribute to results, so accuracy is not wholly autonomous. \\
SLR0629 & \href{https://arxiv.org/pdf/2403.04123}{Exploring LLM-Based Agents for Root Cause Analysis}\par\emph{Retrieval and ReAct comparison} & Private historical records. The text describes a 100-case qualitative sample but table denominators sum to 97. \\
SLR0635 & \href{https://arxiv.org/pdf/2602.06223}{Scaling Mobile Chaos Testing with AI-Driven Test Execution}\par\emph{Industrial diagnostic workflow} & Proprietary system. Diagnostic sample size and independent label protocol are insufficiently specified. \\
SLR0669 & \href{https://arxiv.org/pdf/2510.04711}{Rethinking the Evaluation of Microservice RCA with a Fault Propagation-Aware Benchmark}\par\emph{Diagnostic benchmark} & No LLM evaluated. Fault distribution is imbalanced, impact filtering omits subtle failures and supervised methods use a test subset while other methods use the full dataset. \\
SLR0680 & \href{https://cdn.amazon.science/9d/ff/bc236d704789bdfcfe521bd55614/scipub-approval152129-45181258-expansioncontraction-a-multiagent-graph-traversal-pattern-for-compound-ai-systems.pdf}{Expansion-Contraction: A Multi-Agent Graph Traversal Pattern for Compound AI Systems}\par\emph{Graph-structured multi-agent diagnosis} & Microservice evidence is constructed, not a live telemetry study. Main production results are outside the domain. \\
SLR0685 & \href{https://arxiv.org/pdf/2412.17015}{RCAEval: A Benchmark for Root Cause Analysis of Microservice Systems with Telemetry Data}\par\emph{Diagnostic benchmark} & Injected benchmark applications and a limited reported experiment. Reconstructed baselines and modality adaptations require care. \\
GS1-03 & \href{https://arxiv.org/pdf/2603.00468}{Cloud-OpsBench: A reproducible benchmark for agentic root cause analysis in cloud systems}\par\emph{Diagnostic benchmark} & Frozen snapshots control the observation space but do not reproduce a changing live environment. Predefined tool responses and evidence graphs constrain acceptable investigations. \\
GS1-04 & \href{https://link.springer.com/content/pdf/10.1007/s40747-025-02096-0.pdf?error=cookies_not_supported\&code=1ef4fe7b-5190-45d9-80b8-240de4d8c9c0}{Leveraging multi-agent framework for root cause analysis}\par\emph{Specialised multi-agent workflow} & Small Nezha dataset of 101 cases before splitting. Adapted baselines, unclear uncertainty and resource matching. \\
GS1-06 & \href{https://arxiv.org/pdf/2602.09937}{Why do AI agents systematically fail at cloud root cause analysis?}\par\emph{Empirical evaluation} & One agent framework. Mitigations use smaller Bank subsets (70 prompt tasks and 41 communication cases). \\
GS1-08 & \href{https://arxiv.org/pdf/2605.03505}{LATS-RCA: Language Agent Tree Search for Root Cause Analysis in Microservices}\par\emph{Tree-search multi-agent diagnosis} & One small API-error dataset, pre-identified anomalies, no healthy cases. More search receives substantially more inference resources. \\
GS1-31 & \href{https://ojs.aaai.org/index.php/AAAI/article/view/36991/40953}{Rcaflow: A workflow-informed hierarchical planning multi-agent system for root cause analysis}\par\emph{Workflow-guided multi-agent diagnosis} & Adapted baselines, unclear repeated-run uncertainty, and path length restricted to completed diagnoses. Ablation table has a minor Market value inconsistency. \\
GS1-32 & \href{https://www.techscience.com/cmc/v88n1/67280/html}{Graph-Augmented Multi-Agent Robust Root Cause Analysis in AIOps}\par\emph{Graph-guided multi-agent diagnosis} & Validation-set tuning is described without a clear independent test split. Narrow baseline selection and uncertain repetition units limit confidence. \\
GS1-44 & \href{https://aclanthology.org/2024.findings-emnlp.232.pdf}{mABC: multi-Agent Blockchain-Inspired Collaboration for root cause analysis in micro-services architecture}\par\emph{Multi-agent diagnosis} & Some prose and table values disagree. Exact diagnostic-case denominator is not clear from call-chain counts. \\
GS1-48 & \href{https://arxiv.org/pdf/2601.04709}{Bridging temporal and textual modalities: A multimodal framework for automated cloud failure root cause analysis}\par\emph{Hybrid time-series and retrieval diagnosis} & Forecasting extends short sequences and may change the evidence. Pairwise preference is not verified cause accuracy. \\
GS2-03 & \href{https://nkcs.iops.ai/wp-content/uploads/2026/03/LocaleXpert.pdf}{LLM-Enhanced Failure Localization in Microservices: Integrating Multi-Modal Data and Expert Interpretation}\par\emph{Hybrid modality experts} & Private data and partially sampled human review. Text similarity and preference do not prove causal correctness. \\
GS2-07 & \href{https://arxiv.org/pdf/2607.04623}{Can LLMs Really Recover Microservice Failures? A Recovery-Aware Evaluation of Diagnosis-to-Action Reasoning}\par\emph{Diagnosis-to-action benchmark} & Incomplete cases are excluded by the quality gate. Fixed mappings give RCA-only methods recovery actions they do not generate themselves. \\
GS2-22 & \href{https://arxiv.org/pdf/2607.08529}{Log-Insight: Automating Microservice Incident Diagnosis via Neuro-Symbolic Log Analysis}\par\emph{Statistical preprocessing with LLM synthesis} & Only 11 distinct incidents despite 110 runs. Baseline comparison covers three cases. \\
GS2-28 & \href{https://arxiv.org/pdf/2512.19697}{Automated fault detection in 5G core networks using large language models}\par\emph{Fixed LLM diagnostic classifier} & Small single-testbed dataset. Prompt explicitly supplies fault-to-symptom rules and assumes one fault. \\
GS2-30 & \href{https://www.mdpi.com/1424-8220/26/17/5571/pdf}{A Knowledge-Enhanced Iterative Reasoning Framework for Accurate and Traceable Fault Diagnosis in Distributed Service Systems}\par\emph{Knowledge-graph and LLM hybrid} & One small, closed fault taxonomy with strong class-specific rules. Perfect results on this split do not show generalisation to unknown failures. \\
GS2-37 & \href{https://arxiv.org/pdf/2602.08804}{Root Cause Analysis Method Based on Large Language Models with Residual Connection Structures}\par\emph{Telemetry filtering with LLM diagnosis} & Successive development iterations may overfit the benchmark. No held-out development protocol, repeated uncertainty or strong external baselines is established in the evaluated section. \\
GS2-41 & \href{https://arxiv.org/pdf/2608.21310}{Beyond Fault Localization: A Trajectory-Level Study of LLM Agents for Microservice Root Cause Analysis}\par\emph{Trajectory evaluation and verification intervention} & Preprint, framework adaptations and limited backbones. Agreement for most annotation tasks concerns binary acceptance of proposed labels, not independent initial construction. \\
GS3-03 & \href{https://arxiv.org/pdf/2512.16381}{A network arena for benchmarking ai agents on network troubleshooting}\par\emph{Network diagnostic benchmark} & Mixed network settings and emulation limit cloud transfer. Boolean-mask accuracy is not necessarily case-level exact match. \\
GS3-04 & \href{https://proceedings.mlsys.org/paper_files/paper/2025/file/d1f9e4a9f109b6e8b75ed362736f22ec-Paper-Conference.pdf}{Aiopslab: A holistic framework to evaluate ai agents for enabling autonomous clouds}\par\emph{Live agent benchmark} & FLASH is a simplified reimplementation. The stated 48-problem evaluation and table-implied denominators do not reconcile cleanly. \\
GS3-09 & \href{https://arxiv.org/pdf/2606.29193}{A Multi-Dataset Benchmark for Evaluating LLM Agents in Microservice Failure Diagnosis}\par\emph{Diagnostic benchmark} & Demo scale, curated recoverable cases and subjective composite weights. Keyword/semantic matches can reward superficial evidence mentions. \\
GS3-11 & \href{https://arxiv.org/pdf/2411.05349}{Enhancing cluster resilience: Llm-agent based autonomous intelligent cluster diagnosis system and evaluation framework}\par\emph{Cluster diagnosis agent and benchmark} & Some explicitly marked configurations expose the full dataset or benchmark-specific plans and cannot support held-out effectiveness. GPU-cluster setting is adjacent to the thesis cloud API. \\
GS3-15 & \href{https://arxiv.org/pdf/2604.23455}{CUJBench: Benchmarking LLM Agents for Browser-to-Backend Failure Diagnosis}\par\emph{Browser-to-backend diagnostic benchmark} & Only 40 of 87 reference labels were checked by a human. Four timeout runs excluded, and denominators differ between tables. \\
GS3-24 & \href{https://arxiv.org/pdf/2605.07161}{SREGym: A Live Benchmark for AI SRE Agents with High-Fidelity Failure Scenarios}\par\emph{Live diagnostic benchmark} & Synthetic Kubernetes settings and LLM-judge dependence. Commercial models and unequal token consumption complicate architecture comparisons. \\
GS3-38 & \href{https://arxiv.org/pdf/2608.27021}{FaulT-Bench: Towards Benchmarking Network Troubleshooting LLM Agents under Unreliable User Tickets}\par\emph{Diagnostic reliability benchmark} & Mostly enterprise/campus network settings, limited direct cloud transfer. LLM judge and composite outcome weighting. \\
GS3-39 & \href{https://arxiv.org/pdf/2607.22165}{DBA-Bench: A Production-Fidelity Benchmark for LLM-Based Database Operations Agents}\par\emph{Database operational benchmark} & Diagnosis excludes inapplicable and unresolved judge outputs, creating a different denominator from outcome scores. Mixed tuning/change tasks are not all diagnostic. \\
GS3-48 & \href{https://proceedings.neurips.cc/paper_files/paper/2025/file/47a6e9e2c3019f13ad94a0f259fe4970-Paper-Conference.pdf}{Stratus: A multi-agent system for autonomous reliability engineering of modern clouds}\par\emph{Multi-agent reliability workflow} & RCA labels can be ambiguous, and localisation is scored as Pass@3. With GPT-4o it fails all three no-fault detection problems. Some mitigation successes exploit nonpersistent injections. \\
SUP-OpenRCA & \href{https://proceedings.iclr.cc/paper_files/paper/2025/file/d29b8d53678015079e1d245c023e49d2-Paper-Conference.pdf}{OpenRCA: Can Large Language Models Locate the Root Cause of Software Failures?}\par\emph{Diagnostic benchmark and controller--executor agent} & Known candidate reasons/components and a failure window are supplied. Oracle sampling uses privileged information. \\
SUP-ITBench & \href{https://tingdai.github.io/files/itbench_ICML25.pdf}{ITBench: Evaluating AI Agents across Diverse Real-World IT Automation Tasks}\par\emph{Live operational benchmark} & Only 11 of 102 total-domain scenarios initially released. Table footnote and prose differ on how trace conditions are aggregated, so a modality effect should not be inferred from the headline score alone. \\
SLR0584 & \href{https://www.vldb.org/pvldb/vol19/p4678-zhu.pdf}{SuperDRI: Graph-Guided Autonomous Troubleshooting for Cloud Databases}\par\emph{Graph-guided multi-agent incident troubleshooting (industrial demo)} & Process agreement and query execution are measured, but diagnostic correctness, a matched baseline and model identity are not established. \\
SLR0696 & \href{https://www.vldb.org/pvldb/vol19/p3928-chen.pdf}{DBAgent: An RL-Based Agent for Autonomous Database Operations and Maintenance}\par\emph{RL-fine-tuned LLM agent for database diagnosis and O\&M (DBAgent)} & Joint cause-and-remedy metric. Overall averages include consultation. Training access differs from prompt-only baselines. \\
GS2-42 & \href{https://sol.sbc.org.br/index.php/semish/article/download/43509/43272}{Cloud Microservices: An Empirical Study of MCP-Orchestrated LLM Agents for Incident Triage}\par\emph{Controlled model/temperature comparison of an MCP tool-using agent} & Three scenarios only. The headline success rate pools fault and healthy cases. False-positive rates require separate interpretation. \\
GS2-48 & \href{https://assets-eu.researchsquare.com/files/rs-10499343/v1_covered_585c57a7-39b6-4e04-a86f-bdcf416ada4e.pdf?c=1790145669}{Agent-MIRUPD: Operationalizing LLM Agents for Microservices Incident Response Under Performance Degradation}\par\emph{Staged lifecycle agent with stage-context propagation and HITL mitigation} & Preprint with small experiments. Published baseline results were reused rather than rerun under common conditions. \\
SLR0563 & \href{https://www.researchgate.net/publication/402677372_Representation-Aware_Root_Cause_Analysis_with_Large_Language_Models_Position_Paper}{Representation-Aware Root Cause Analysis with Large Language Models (Position Paper)}\par\emph{Representation-controlled LLM diagnosis} & Single system and model. TraceRCA is a contextual reference, not a controlled head-to-head baseline. \\
SLR0009 & \href{https://www.researchgate.net/publication/407161871_A_Retrieval-Augmented_Log-Metric-Trace_Twin_for_Backend_Failure_Diagnosis}{A Retrieval-Augmented Log-Metric-Trace Twin for Backend Failure Diagnosis}\par\emph{Structured evidence and historical retrieval} & Synthetic incidents. Explanation ratings and localisation are different outcomes. \\
SLR0048 & \href{https://www.researchgate.net/publication/404197467_CausalRCA_Causal_Graph-Augmented_Retrieval_for_Root_Cause_Analysis_in_Microservice_Systems}{CausalRCA: Causal Graph-Augmented Retrieval for Root Cause Analysis in Microservice Systems}\par\emph{Graph and retrieval guided diagnosis} & One testbed and template-generated knowledge. The confidence score is not established as a calibrated probability. \\
SLR0167 & \href{https://www.researchgate.net/publication/407715364_VerifyRCA_LLM-Guided_but_Evidence-Verified_Root_Cause_Analysis_for_Microservice_Backends}{VerifyRCA: LLM-Guided but Evidence-Verified Root Cause Analysis for Microservice Backends}\par\emph{Hypothesis generation and evidence verification} & The ablation also removes evidence queries. Confidence is a ranking margin, not a calibrated probability. \\
SLR0183 & \href{https://www.researchgate.net/publication/405467638_Evidence-Verified_Root_Cause_Localization_for_Microservice_Incidents_under_Tool_and_Telemetry_Noise}{Evidence-Verified Root Cause Localization for Microservice Incidents under Tool and Telemetry Noise}\par\emph{Evidence verification under noise} & Synthetic perturbations and inconsistent category descriptions limit transfer. Evidence ratings use each method’s best clean run. \\
SLR0266 & \href{https://www.researchgate.net/publication/406279543_Evidence-Verified_LLM_Agents_for_Safe_Backend_Incident_Remediation}{Evidence-Verified LLM Agents for Safe Backend Incident Remediation}\par\emph{Diagnosis followed by an evidence gate} & Wrong-fix denominator includes escalated incidents. Lower harm partly reflects deferral, not more accurate diagnosis. \\
SLR0328 & \href{https://www.researchgate.net/publication/405785833_AgentQE-Bench_A_Reproducible_Evaluation_Framework_for_Agentic_AI_Reliability_Safety_and_ROI_in_CICD_Pipelines}{AgentQE-Bench: A Reproducible Evaluation Framework for Agentic AI Reliability, Safety, and ROI in CI/CD Pipelines}\par\emph{Simulated diagnosis and action benchmark} & Small deterministic simulator. Composite scores obscure diagnostic performance. Check observations for label leakage. \\
SLR0613 & \href{https://arxiv.org/pdf/2309.05833v3}{LM-PACE: Confidence Estimation by Large Language Models for Effective Root Causing of Cloud Incidents}\par\emph{Retrieval-assisted confidence calibration} & Pseudo-label dependence and private data. Reported split sizes sum to 103,308, not the stated 121,308 total. \\
\end{longtable}
\endgroup
